% Generated draft: requires mathematical and source review.
\documentclass[11pt]{article}
\usepackage[margin=25mm]{geometry}
\usepackage{fontspec}
\setmainfont{Latin Modern Roman}
\usepackage{amsmath,amssymb,mathtools,unicode-math}
\setmathfont{Latin Modern Math}
\usepackage{enumitem,needspace,xurl,hyperref}
\setlength{\parindent}{0pt}
\setlength{\parskip}{5pt}
\setlength{\emergencystretch}{3em}
\newcommand{\N}{\mathbb{N}}
\newcommand{\Z}{\mathbb{Z}}
\newcommand{\Q}{\mathbb{Q}}
\newcommand{\R}{\mathbb{R}}
\newcommand{\C}{\mathbb{C}}
\newcommand{\F}{\mathbb{F}}
\newcommand{\A}{\mathbb{A}}
\newcommand{\PP}{\mathbb{P}}
\newcommand{\E}{\mathbb{E}}
\newcommand{\eps}{\varepsilon}
\newcommand{\dd}{\,\mathrm{d}}
\begin{document}
\section*{UnsolvedMath ID: 800673}
\textbf{Source problem code:} catalogue entry
\subsection{Definitions and mathematical statement}
For a topological space \(X\), use the customary real-valued interpretation
\[
C_p(X)=\{f:X\to\R:f\text{ is continuous}\},
\]
with the topology inherited from the product \(\R^X\). Thus a basic neighborhood of \(f\in C_p(X)\) is specified by a finite set \(F\subseteq X\) and \(\varepsilon>0\), and consists of all \(g\in C_p(X)\) such that \(|g(x)-f(x)|<\varepsilon\) for every \(x\in F\).

A space \(Y\) has the Pytkeev property if, whenever \(A\subseteq Y\) and \(y\in\overline{A\setminus\{y\}}\), there is a countable family \(\mathcal P\) of infinite subsets of \(A\setminus\{y\}\) such that every neighborhood \(V\) of \(y\) contains some \(P\in\mathcal P\).

Use the following standard countable-cover convention to express the selection property. An open cover of a space \(Y\) is a family of open subsets whose union is \(Y\), with \(Y\) itself excluded as a member. An open \(\omega\)-cover is such a cover in which every finite subset of \(Y\) is contained in one member. An open \(\gamma\)-cover is a countably infinite such cover \(\mathcal U\) for which
\[
\bigl|\{U\in\mathcal U:y\notin U\}\bigr|<\infty
\qquad(y\in Y).
\]
Write \(\Gamma(Y)\) for the family of these \(\gamma\)-covers. The property \(\mathsf U_{\mathrm{fin}}(\Gamma,\Gamma)\) means that for every sequence \((\mathcal U_k)_{k\geq1}\) in \(\Gamma(Y)\), with no \(\mathcal U_k\) having a finite subcover, there are finite subfamilies \(\mathcal F_k\subseteq\mathcal U_k\) such that
\[
\left\{\bigcup\mathcal F_k:k\geq1\right\}\in\Gamma(Y).
\]
Finite powers \(X^n\), for positive integers \(n\), carry the product topology.

The question attributed in the supplied record to \([\mathrm{SakaiRez},\ \text{Question }2.9]\), and reproduced as Problem 46 in Tsaban's list, is whether
\[
C_p(X)\text{ has the Pytkeev property}
\quad\Longrightarrow\quad
(\forall n\geq1)\;X^n\models\mathsf U_{\mathrm{fin}}(\Gamma,\Gamma).
\]
Here \(X\) must range over the class intended by the source. The supplied excerpt does not specify that class or its separation assumptions, so the scope of this quantifier remains unresolved. The real-valued interpretation of \(C_p(X)\) and the countable-cover conventions above are explicit standard interpretations requiring confirmation against the complete source; they are not additional hypotheses supplied by the record.

The excerpt also supplies the contextual implication that, if every finite power of \(X\) satisfies \(\mathsf U_{\mathrm{fin}}(\Gamma,\Gamma)\), then every open \(\omega\)-cover of \(X\) is \(\omega\)-shrinkable, attributed to \([\mathrm{SakaiRez}]\). It says that a positive answer to the displayed question would therefore suffice for an antecedent objective. Neither the definition of \(\omega\)-shrinkability nor that antecedent objective is supplied, and both remain unresolved here.
\subsection{Short English statement}
The supplied question asks whether the Pytkeev property of the pointwise topology on the continuous real-valued functions on \(X\) forces every positive finite power of \(X\) to satisfy \(\mathsf U_{\mathrm{fin}}(\Gamma,\Gamma)\). The intended class of spaces and the source's cover conventions require confirmation. The registry labels the problem open; this is a supplied status, not an independently verified literature assessment.
\subsection{Sources}
\begin{itemize}
\item Supplied local UnsolvedMath registry record: permanent ID \(800673\), exact source problem code \texttt{catalogue entry}; title: Selection Principles Milestone Problems --- Problem 46 ([SakaiRez, Question 2.9]). The record lists Tsaban as proposer, the year as 2003, and the status as open.
\item Supplied source reference: Tsaban, \emph{Selection Principles in Mathematics: A Milestone of Open Problems} (2003), Problem 46, \url{https://arxiv.org/abs/math/0312182}. The record reports access on 2026-10-16 and extraction by \texttt{arxiv\_latex}. These are supplied provenance details; no source visit or independent verification is claimed.
\item Attribution retained from the excerpt: \([\mathrm{SakaiRez},\ \text{Question }2.9]\). The contextual assertion concerning \(\omega\)-shrinkable covers is also attributed there to \([\mathrm{SakaiRez}]\). The supplied record does not provide a full bibliographic entry or a separate URL for this reference.
\end{itemize}
\end{document}
